% =============================================================================
% Section 2: Background and Threat Model
% (sec_threat_model.tex content merged here)
% =============================================================================

\section{Background and Threat Model}
\label{sec:background}

\subsection{Membership Inference Attacks}

Let $D = \{(x_i, y_i)\}_{i=1}^{n}$ be a training dataset drawn i.i.d.\
from an unknown distribution $\mathcal{P}$ over $\mathcal{X} \times \mathcal{Y}$.
A learner trains a classifier $f_\theta : \mathcal{X} \to \Delta(\mathcal{Y})$
by minimizing empirical risk on $D$.
A membership inference attack (MIA) is an algorithm
$\mathcal{A}$ that, given a target sample $(x, y)$ and query access to
$f_\theta$, outputs $1$ if $(x,y) \in D$ and $0$ otherwise.

The MIA advantage of $\mathcal{A}$ is:
\begin{equation}
  \mathrm{Adv}(\mathcal{A}, f_\theta, D)
  =
  \bigl|
    \Pr[\mathcal{A}=1 \mid (x,y)\in D]
    -
    \Pr[\mathcal{A}=1 \mid (x,y)\notin D]
  \bigr|.
  \label{eq:adv}
\end{equation}
We denote the dataset-level risk as the expected advantage over samples:
\begin{equation}
  \mathrm{Risk}(D) = \sup_{\mathcal{A}} \;\mathbb{E}_{(x,y)\sim\mathcal{P}}\,
  \mathrm{Adv}(\mathcal{A}, f_\theta, D).
  \label{eq:risk}
\end{equation}
In practice, $\mathrm{Risk}(D)$ is estimated via AUC of the best
empirical attack over a grid of attacks and model families
(Section~\ref{sec:setup}).

\subsection{Differential Privacy}

An algorithm $\mathcal{M}$ is $(\varepsilon,\delta)$-differentially
private~\cite{dwork2006calibrating} if for all adjacent datasets $D, D'$
differing in one record and all outputs $S$:
$\Pr[\mathcal{M}(D)\in S] \leq e^\varepsilon \Pr[\mathcal{M}(D')\in S] + \delta$.
Yeom et al.~\cite{yeom2018privacy} showed that
$\mathrm{Adv}(\mathcal{A}^*) \leq e^\varepsilon - 1 + \delta$
for $(\varepsilon,\delta)$-DP mechanisms.
Our work shows that this bound, while tight in the worst case, does not
account for how much DPRI amplifies or attenuates the practical risk
(Section~\ref{sec:finding3}).

\subsection{Threat Model}
\label{sec:threat}

We consider two principals: a data holder who trains a model on
$D$ and may later publish it, and an adversary who infers membership.

\paragraph{Adversary goal.}
Given $(x,y)$ and oracle access to $f_\theta$, determine whether
$(x,y) \in D$.

\paragraph{Adversary capability.}
Black-box access: the adversary submits inputs and observes
predicted class probabilities, but cannot inspect model weights.
This is the standard setting for deployed APIs and MLaaS platforms.
White-box results are consistent and reported in Appendix~\ref{app:whitebox}.

\paragraph{Adversary knowledge.}
The adversary knows the data domain distribution $\mathcal{P}$
(e.g., that the dataset contains medical records of a specific demographic)
but does not know the exact set $D$.

\paragraph{Defender goal.}
Before training, compute $\mathrm{DPRI}(D)$ from raw data alone, with no
model required, to estimate $\mathrm{Risk}(D)$ and inform DP budget
selection and release decisions.

\paragraph{Dual use.}
DPRI also has an offensive application: an adversary with access to
candidate target datasets can compute DPRI pre-attack and prioritize
high-risk datasets to maximize attack return on query budget.
We quantify this in Section~\ref{sec:implications} and discuss mitigations.

\paragraph{Scope.}
Centralized training, passive adversary (no data poisoning).
Extensions to federated learning are discussed in Section~\ref{sec:implications}.
